\documentclass{amsart}
% For dealing with references we use the comment environment
\usepackage{verbatim}
\newenvironment{reference}{\comment}{\endcomment}
%\newenvironment{reference}{}{}
\newenvironment{slogan}{\comment}{\endcomment}
\newenvironment{history}{\comment}{\endcomment}

% For commutative diagrams we use Xy-pic
\usepackage[all]{xy}

% We use 2cell for 2-commutative diagrams.
\xyoption{2cell}
\UseAllTwocells

% We use multicol for the list of chapters between chapters
\usepackage{multicol}

% This is generally recommended for better output
\usepackage{lmodern}
\usepackage[T1]{fontenc}

% For cross-file-references

% Package for hypertext links:
\usepackage{hyperref}

% For any local file, say "hello.tex" you want to link to please

% Theorem environments.
%
\theoremstyle{plain}
\newtheorem{theorem}[subsection]{Theorem}
\newtheorem{proposition}[subsection]{Proposition}
\newtheorem{lemma}[subsection]{Lemma}

\theoremstyle{definition}
\newtheorem{definition}[subsection]{Definition}
\newtheorem{example}[subsection]{Example}
\newtheorem{exercise}[subsection]{Exercise}
\newtheorem{situation}[subsection]{Situation}

\theoremstyle{remark}
\newtheorem{remark}[subsection]{Remark}
\newtheorem{remarks}[subsection]{Remarks}

\numberwithin{equation}{subsection}

% Macros
%
\def\lim{\mathop{\mathrm{lim}}\nolimits}
\def\colim{\mathop{\mathrm{colim}}\nolimits}
\def\Spec{\mathop{\mathrm{Spec}}}
\def\Hom{\mathop{\mathrm{Hom}}\nolimits}
\def\Ext{\mathop{\mathrm{Ext}}\nolimits}
\def\SheafHom{\mathop{\mathcal{H}\!\mathit{om}}\nolimits}
\def\SheafExt{\mathop{\mathcal{E}\!\mathit{xt}}\nolimits}
\def\Sch{\mathit{Sch}}
\def\Mor{\mathop{\mathrm{Mor}}\nolimits}
\def\Ob{\mathop{\mathrm{Ob}}\nolimits}
\def\Sh{\mathop{\mathit{Sh}}\nolimits}
\def\NL{\mathop{N\!L}\nolimits}
\def\CH{\mathop{\mathrm{CH}}\nolimits}
\def\proetale{{pro\text{-}\acute{e}tale}}
\def\etale{{\acute{e}tale}}
\def\QCoh{\mathit{QCoh}}
\def\Ker{\mathop{\mathrm{Ker}}}
\def\Im{\mathop{\mathrm{Im}}}
\def\Coker{\mathop{\mathrm{Coker}}}
\def\Coim{\mathop{\mathrm{Coim}}}

% Boxtimes
%
\DeclareMathSymbol{\boxtimes}{\mathbin}{AMSa}{"02}

%
% Macros for moduli stacks/spaces
%
\def\QCohstack{\mathcal{QC}\!\mathit{oh}}
\def\Cohstack{\mathcal{C}\!\mathit{oh}}
\def\Spacesstack{\mathcal{S}\!\mathit{paces}}
\def\Quotfunctor{\mathrm{Quot}}
\def\Hilbfunctor{\mathrm{Hilb}}
\def\Curvesstack{\mathcal{C}\!\mathit{urves}}
\def\Polarizedstack{\mathcal{P}\!\mathit{olarized}}
\def\Complexesstack{\mathcal{C}\!\mathit{omplexes}}
% \Pic is the operator that assigns to X its picard group, usage \Pic(X)
% \Picardstack_{X/B} denotes the Picard stack of X over B
% \Picardfunctor_{X/B} denotes the Picard functor of X over B
\def\Pic{\mathop{\mathrm{Pic}}\nolimits}
\def\Picardstack{\mathcal{P}\!\mathit{ic}}
\def\Picardfunctor{\mathrm{Pic}}
\def\Deformationcategory{\mathcal{D}\!\mathit{ef}}

\setlength{\emergencystretch}{3em}
\begin{document}
\title[Stacks Project, Tag 0A6W]{243 --- Greenlees-May identities relating derived completion and local cohomology}
\maketitle
\noindent Copyright (C) 2005--2025 Johan de Jong. Stacks Project authors. Source: \href{https://stacks.math.columbia.edu/tag/0A6W}{Tag 0A6W}.

\noindent Editorial difficulty note (not author text). Difficulty H2: Both identities are proved by constructing and filtering the two different mapping cones. The graded pieces are derived Hom from localizations or localized derived tensor products, and distinct vanishing mechanisms annihilate them. The task requires compatible control of derived completion and torsion on arbitrary complexes..

\noindent Editorial provenance note (not author text). History: excerpt from revision \texttt{a0444\allowbreak 6e57e\allowbreak c1fbc\allowbreak 252a8\allowbreak 71afc\allowbreak ec775\allowbreak 2fb28\allowbreak 07b14}, dualizing.tex, lines 1952--2012. Reference presentation was adapted; original-author context and core support blocks were retained or added. The mathematical wording is unchanged. Licensed under GNU FDL 1.2 or later; no invariant sections or cover texts. See ../licenses/COPYING. Context blocks reproduce author definitions and supporting results; they are not additional counted problems.

\begin{lemma}
\label{lemma-local-cohomology-adjoint}
Let $A$ be a ring and let $I \subset A$ be a finitely generated ideal.
There exists a right adjoint $R\Gamma_Z$ (\href{https://stacks.math.columbia.edu/tag/0A6Q}{equation local cohomology (Tag 0A6Q)})
to the inclusion functor $D_{I^\infty\text{-torsion}}(A) \to D(A)$.
In fact, if $I$ is generated by $f_1, \ldots, f_r \in A$, then we have
$$
R\Gamma_Z(K) =
(A \to \prod\nolimits_{i_0} A_{f_{i_0}} \to
\prod\nolimits_{i_0 < i_1} A_{f_{i_0}f_{i_1}}
\to \ldots \to A_{f_1\ldots f_r}) \otimes_A^\mathbf{L} K
$$
functorially in $K \in D(A)$.
\end{lemma}

\begin{lemma}
\label{more-algebra-lemma-derived-completion}
\begin{slogan}
Derived completions along finitely generated ideals exist, and can
be computed by a {\v C}ech procedure.
\end{slogan}
Let $I$ be a finitely generated ideal of a ring $A$.
The inclusion functor $D_{comp}(A, I) \to D(A)$ has a
left adjoint, i.e., given any object $K$ of $D(A)$ there
exists a map $K \to K^\wedge$ of $K$ into a derived complete
object of $D(A)$ such that the map
$$
\Hom_{D(A)}(K^\wedge, E) \longrightarrow \Hom_{D(A)}(K, E)
$$
is bijective whenever $E$ is a derived complete object of $D(A)$.
In fact, if $I$ is generated by $f_1, \ldots, f_r \in A$, then we have
$$
K^\wedge = R\Hom\left((A \to \prod\nolimits_{i_0} A_{f_{i_0}} \to
\prod\nolimits_{i_0 < i_1} A_{f_{i_0}f_{i_1}}
\to \ldots \to A_{f_1\ldots f_r}), K\right)
$$
functorially in $K$.
\end{lemma}

\begin{lemma}
\label{lemma-complete-and-local}
\begin{slogan}
Results of this nature are sometimes referred to as Greenlees-May duality.
\end{slogan}
Let $A$ be a ring and let $I$ be a finitely generated ideal.
Let $R\Gamma_Z$ be as in Lemma \ref{lemma-local-cohomology-adjoint}.
Let ${\ }^\wedge$ denote derived completion as in
More on Algebra, Lemma \ref{more-algebra-lemma-derived-completion}.
For an object $K$ in $D(A)$ we have
$$
R\Gamma_Z(K^\wedge) = R\Gamma_Z(K)
\quad\text{and}\quad
(R\Gamma_Z(K))^\wedge = K^\wedge
$$
in $D(A)$.
\end{lemma}

\begin{proof}
Choose $f_1, \ldots, f_r \in A$ generating $I$. Recall that
$$
K^\wedge = R\Hom_A\left((A \to \prod A_{f_{i_0}}
\to \prod A_{f_{i_0i_1}} \to \ldots \to A_{f_1 \ldots f_r}), K\right)
$$
by More on Algebra, Lemma \ref{more-algebra-lemma-derived-completion}.
Hence the cone $C = \text{Cone}(K \to K^\wedge)$
is given by
$$
R\Hom_A\left((\prod A_{f_{i_0}}
\to \prod A_{f_{i_0i_1}} \to \ldots \to A_{f_1 \ldots f_r}), K\right)
$$
which can be represented by a complex endowed with a finite filtration
whose successive quotients are isomorphic to
$$
R\Hom_A(A_{f_{i_0} \ldots f_{i_p}}, K), \quad p > 0
$$
These complexes vanish on applying $R\Gamma_Z$, see
Lemma \href{https://stacks.math.columbia.edu/tag/0A6S}{lemma local cohomology vanishes (Tag 0A6S)}. Applying $R\Gamma_Z$
to the distinguished triangle $K \to K^\wedge \to C \to K[1]$
we see that the first formula of the lemma is correct.

\medskip\noindent
Recall that
$$
R\Gamma_Z(K) =
K \otimes^\mathbf{L} (A \to \prod A_{f_{i_0}}
\to \prod A_{f_{i_0i_1}} \to \ldots \to A_{f_1 \ldots f_r})
$$
by Lemma \ref{lemma-local-cohomology-adjoint}.
Hence the cone $C = \text{Cone}(R\Gamma_Z(K) \to K)$
can be represented by a complex endowed with a finite filtration
whose successive quotients are isomorphic to
$$
K \otimes_A A_{f_{i_0} \ldots f_{i_p}}, \quad p > 0
$$
These complexes vanish on applying ${\ }^\wedge$, see
More on Algebra, Lemma \href{https://stacks.math.columbia.edu/tag/0A6D}{lemma derived completion vanishes (Tag 0A6D)}.
Applying derived completion to the distinguished triangle
$R\Gamma_Z(K) \to K \to C \to R\Gamma_Z(K)[1]$
we see that the second formula of the lemma is correct.
\end{proof}
\end{document}
